\begin{table*}[ht!]
\small
\caption{Examples of model predictions. Gold leaf is highlighted in \textbf{bold} in Table of Contents}
\label{table:pred_examples}
\resizebox{\textwidth}{!}{%
\begin{tabular}{l|l}
\multicolumn{2}{c}{\textbf{Example 1:}} \\
\toprule
\multicolumn{2}{l}{\textbf{Book}: \href{https://open.umn.edu/opentextbooks/textbooks/1621}{The Whole Child: Development in the Early Years}}\\
\multicolumn{2}{l}{\textbf{Query}: How do preschoolers react when caregivers and teachers belittle their autonomous actions?}\\
\midrule
\textbf{Table of Contents} & \textbf{Content}\\
$\cdots$\\
2 Chapter Two: Theorists and Theories of Development & Initiative vs. Guilt (Preschool Years)\\
leaf 2.1 Theories of Development & The development of courage and independence ...\\
2.2 Psychosocial Theory & ...\\
leaf 2.2.1 Trust vs. Mistrust (Infancy)& ...\\
leaf 2.2.2 Autonomy vs. Shame/Doubt (Toddlerhood) &...\\
\textbf{leaf 2.2.3 Initiative vs. Guilt (Preschool Years)}\\
leaf 2.2.4 Industry vs. Inferiority (Elementary Years) & If caregivers and preschool teachers encourage and support\\
leaf 2.2.5 Identity vs. Role-Confusion (Adolescence) & children’s efforts while also helping them make realistic \\
leaf 2.2.6 Intimacy vs. Isolation (Early Adulthood) & and appropriate choices, children develop a healthy sense of initiative in \\
leaf 2.2.7 Generativity vs. Stagnation (Middle Adulthood)& planning and undertaking activities If, instead, adults
discourage the\\
leaf 2.2.8 Integrity vs. Despair (Older Adulthood) & pursuit of independent activities or dismiss them as silly and \\
$\cdots$ & bothersome, children develop guilt about their needs and desires.\\
3 Chapter Three: Domains in Development\\
\midrule
\multicolumn{2}{l}{\textbf{Mistral}: 8.3.3 Moral Development, leaf 7.3.2 Hitting/Scratching and Temper Tantrums, \textbf{leaf 2.2.3 Initiative vs. Guilt (Preschool Years)}}\\
\multicolumn{2}{l}{\textbf{BM25}: 8.3.2 Moral Development}\\
\multicolumn{2}{l}{\textbf{DPR: }2.2.2 Autonomy vs. Shame/Doubt (Toddlerhood)}\\
\multicolumn{2}{l}{\textbf{DSI}: 2.2.2 Autonomy vs. Shame/Doubt (Preschool Years)}\\
\multicolumn{2}{l}{\textbf{{\sys}}: \textbf{2.2.3 Initiative vs. Guilt (Preschool Years)}}\\
\midrule
\bottomrule
\multicolumn{2}{c}{\textbf{Example 2:}} \\
\toprule
\multicolumn{2}{l}{\textbf{Book}: \href{https://open.umn.edu/opentextbooks/textbooks/1013}{Nursing Fundamentals}}\\
\multicolumn{2}{l}{\textbf{Query}: In what way does a chronic illness affect an elderly person's ability to perform daily activities?}\\
\midrule
\textbf{Table of Contents} & \textbf{Content}\\
$\cdots$\\
leaf 30.5 Spiritual Care of Self & Applying the Nursing Process\\
leaf 30.6 Putting It All Together& ...\\
31 Care of the Older Adult & It is also important to consider the impact of chronic disease\\
leaf 31.1 Care of the Older Adult Introduction & on
their ability to function and complete Activities of Daily Living\\
leaf 31.2 Basic Concepts &(ADLs). older adults who are able to perform ADLs without assistance\\
leaf \textbf{31.3 Applying the Nursing Process} & consider themselves healthy.\\
\midrule
\multicolumn{2}{l}{\textbf{Mistral}: \textbf{31.3 Applying the Nursing Process}}\\
\multicolumn{2}{l}{\textbf{BM25}: 11.2 Sensory Impairments Basic Concepts}\\
\multicolumn{2}{l}{\textbf{DPR: } 31.2 Basic Concepts}\\
\multicolumn{2}{l}{\textbf{DSI}: 31.2 Basic Concepts}\\
\multicolumn{2}{l}{\textbf{{\sys}}: \textbf{31.3 Applying the Nursing Process}}\\
\midrule
\bottomrule
\multicolumn{2}{c}{\textbf{Example 3:}} \\
\toprule
\multicolumn{2}{l}{\textbf{Book}: \href{https://open.umn.edu/opentextbooks/textbooks/1435}{Introduction to Comparative Government and Politics}}\\
\multicolumn{2}{l}{\textbf{Query}: What are irregular armed organizations and how are they used by states?}\\
\midrule
\textbf{Table of Contents} & \textbf{Content}\\
9 Chapter 5: Non-Democracies and Democratic Backsliding\\
$\cdots$ & Another powerful instrument of repression are paramilitaries.\\
9.2 Strategies for staying in power & These refer to groups with access 
to military-grade weapons and training\\
leaf 9.2.1 \textbf{Institutional channels}& yet they are not part of the national military.\\
leaf 9.2.2 Cultural and ideological controls & They are “irregular armed organizations that carry out acts of violence\\
9.3 Varieties of non-democracy & against civilians on behalf of a state,” \\
$\cdots$\\
\midrule
\multicolumn{2}{l}{\textbf{Mistral}: 9.3.5 Illiberal and hybrid regimes}\\
\multicolumn{2}{l}{\textbf{BM25}: 15.3.1 Insurgencies/Civil Wars}\\
\multicolumn{2}{l}{\textbf{DPR}: 15.2.2 External State-Sponsored Political Violence (State-Sponsored Terrorism)}\\
\multicolumn{2}{l}{\textbf{DSI}: 15.2.1 Internal State-Sponsored Political Violence (Government Terrorism)}\\
\multicolumn{2}{l}{\textbf{{\sys}}: \textbf{9.2.1 Institutional channels}}\\
\end{tabular}
}
\end{table*}


\begin{table*}[tb]
\centering
\small
\caption{Recall@1, Recall@3 and nDCG@3 for {\ben}; Haiku is Claude haiku 4.5}
\label{table:results_r1}
\resizebox{\textwidth}{!}{%
\begin{tabular}{lllllll|l}
\toprule
 & \textbf{Education} & \textbf{Finance} & \textbf{Law} & \textbf{Med} & \textbf{NatSci} & \textbf{SocSci} & \textbf{Avg}\\
 \midrule
Mistral & 14.8/17.8/16.6 & 16.4/19.4/18.2 & 13.6/15.7/14.9 & 13.7/17.0/15.7 & 13.1/15.3/14.4 & 10.9/13.4/12.4 & 13.8/16.4/15.4\\
BM25 & 58.7/75.8/68.8 & 55.0/76.8/67.8 & 62.0/77.7/71.3 & 59.6/78.7/70.8 & 58.7/78.1/70.1 & 62.8/78.3/72.0 & 59.5/77.6/70.1\\
Haiku & 59.8/80.4/72.0 & 37.7/56.3/48.4 & 41.1/60.6/52.3 & 46.1/65.0/57.2& 54.6/73.7/65.7 & 34.0/48.8/42.6 & 45.6/64.1/56.4 \\
DPR & 71.5/86.8/80.6 & 66.4/85.5/77.7 & 69.6/84.1/78.2 & 72.7/88.8/82.2 & 63.0/84.1/75.7 & 68.8/83.1/77.3 & 68.7/85.4/78.6\\
DSI & 76.2/84.1/80.9 & 78.1/87.8/83.9 & 73.3/82.9/79.0 & 82.1/89.9/86.8 & 76.2/84.7/81.3 & 75.4/82.2/79.4 & 76.9/85.3/81.9\\
{\sys} & \textbf{83.3}/\textbf{91.2}/\textbf{88.0} & \textbf{82.8}/\textbf{91.4}/\textbf{88.0} & \textbf{80.8}/\textbf{89.4}/\textbf{85.9} & \textbf{86.1}/\textbf{93.2}/\textbf{90.3} & \textbf{80.6}/\textbf{90.4}/\textbf{86.4} & \textbf{81.8}/\textbf{89.0}/\textbf{86.1} & \textbf{82.6}/\textbf{90.8}/\textbf{87.5}\\
\bottomrule
\end{tabular}
}
\end{table*}


% \begin{table*}[tb]
% \centering
% \small
% \caption{Recall@1, Recall@3 and NDCG@3 for {\ben}}
% \label{table:results_r1}
% \begin{tabular}{lllllll|l}
% \toprule
%  & \textbf{Edu} & \textbf{Fin} & \textbf{Law} & \textbf{Med} & \textbf{Nat} & \textbf{Soc} & \textbf{Avg}\\
% % \hline
% \multicolumn{8}{c}{Recall @1}\\
% \midrule
% Mistral & 0.148 & 0.164 & 0.136 & 0.137 & 0.131 & 0.109 & 0.138\\
% BM25 & 0.587 & 0.55 & 0.62 & 0.596 & 0.587 & 0.628 & 0.595\\
% DPR & 0.715 & 0.664 & 0.696 &0.727 & 0.63 & 0.688 & 0.687	\\
% DSI & 0.762 & 0.781 & 0.733 & 0.821 & 0.762 & 0.754 & 0.769\\
% {\sys} & 0.833 & 0.828 & 0.808 & 0.861 & 0.806 & 0.818 & \textbf{0.826}\\
% \multicolumn{8}{c}{Recall @3}\\
% \midrule
% Mistral & 0.178 & 0.194 & 0.157 & 0.17 & 0.153 & 0.134 &0.164\\ 
% BM25 & 0.758 & 0.768 & 0.777 & 0.787 & 0.781 & 0.783 & 0.776\\
% DPR & 0.868 & 0.855 & 0.841 & 0.888 & 0.841& 0.831 & 0.854\\
% DSI & 0.841 & 0.878 & 0.829 & 0.899 & 0.847 & 0.822 & 0.853\\
% {\sys} & 0.912 & 0.914 & 0.894 & 0.932 & 0.904 & 0.89 & \textbf{0.908}\\
% \multicolumn{8}{c}{nDCG@3}\\
% \midrule
% Mistral & 0.166 & 0.182 & 0.149 & 0.157 & 0.144 & 0.124 & 0.154\\
% BM25 & 0.688 & 0.678 & 0.713 & 0.708 & 0.701 & 0.72 & 0.701\\
% DPR & 0.806 & 0.777 & 0.782 & 0.822 & 0.757 & 0.773 & 	0.786\\
% DSI & 0.809 & 0.839 & 0.79 & 0.868 & 0.813 & 0.794 &  0.819	\\
% {\sys} & 0.88 & 0.88 & 0.859 & 0.903 & 0.864 & 0.861 & \textbf{0.875}\\
% \bottomrule
% \end{tabular}
% \end{table*}


% \begin{table*}[tb]
% \centering
% \caption{Recall@1 for {\ben}}
% \label{table:results_r1}
% \begin{tabular}{lllllll|l}
% \toprule
%  & \textbf{Edu} & \textbf{Fin} & \textbf{Law} & \textbf{Med} & \textbf{Nat} & \textbf{Soc} & \textbf{Avg}\\
% \hline
% Mistral & 0.148 & 0.164 & 0.136 & 0.137 & 0.131 & 0.109 & 0.138\\
% BM25 & 0.587 & 0.55 & 0.62 & 0.596 & 0.587 & 0.628 & 0.595\\
% DPR & 0.715 & 0.664 & 0.696 &0.727 & 0.63 & 0.688 & 0.687	\\
% DSI & 0.762 & 0.781 & 0.733 & 0.821 & 0.762 & 0.754 & 0.769\\
% {\sys} & 0.833 & 0.828 & 0.808 & 0.861 & 0.806 & 0.818 & \textbf{0.826}\\
% \bottomrule
% \end{tabular}
% \end{table*}

% \begin{table*}[tb]
% \centering
% \caption{Recall@3 for {\ben}}
% \label{table:results_r3}
% \begin{tabular}{lllllll|l}
% \toprule
%  & \textbf{Edu} & \textbf{Fin} & \textbf{Law} & \textbf{Med} & \textbf{Nat} & \textbf{Soc} & \textbf{Avg}\\
%  % & \textbf{Education} & \textbf{Finance} & \textbf{Law} & \textbf{Medicine} & \textbf{NatSci} & \textbf{SocSci} & \textbf{Avg}\\
% \hline
% Mistral & 0.178 & 0.194 & 0.157 & 0.17 & 0.153 & 0.134 &0.164\\ 
% BM25 & 0.758 & 0.768 & 0.777 & 0.787 & 0.781 & 0.783 & 0.776\\
% DPR & 0.868 & 0.855 & 0.841 & 0.888 & 0.841& 0.831 & 0.854\\
% DSI & 0.841 & 0.878 & 0.829 & 0.899 & 0.847 & 0.822 & 0.853\\
% {\sys} & 0.912 & 0.914 & 0.894 & 0.932 & 0.904 & 0.89 & \textbf{0.908}\\
% \bottomrule
% \end{tabular}
% \end{table*}

% \begin{table*}[tb]
% \centering
% \caption{NDCG@3 for {\ben}}
% \label{table:ndcg_r3}
% \begin{tabular}{lllllll|l}
% \toprule
%  & \textbf{Edu} & \textbf{Fin} & \textbf{Law} & \textbf{Med} & \textbf{Nat} & \textbf{Soc} & \textbf{Avg}\\
% \hline
% Mistral & 0.166 & 0.182 & 0.149 & 0.157 & 0.144 & 0.124 & 0.154\\
% BM25 & 0.688 & 0.678 & 0.713 & 0.708 & 0.701 & 0.72 & 0.701\\
% DPR & 0.806 & 0.777 & 0.782 & 0.822 & 0.757 & 0.773 & 	0.786\\
% DSI & 0.809 & 0.839 & 0.79 & 0.868 & 0.813 & 0.794 &  0.819	\\
% {\sys} & 0.88 & 0.88 & 0.859 & 0.903 & 0.864 & 0.861 & \textbf{0.875}\\
% \bottomrule
% \end{tabular}
% \end{table*}
Table~\ref{table:results_r1} compares {\sys} with all the baselines listed in Section~\ref{sec:baselines} and addresses the following two research questions:
\begin{itemize}
    \item \textbf{RQ1}: Does {\toc} based training makes {\sys} more accurate as a retriever?
    \item \textbf{RQ2}: Can we finetune an LLM to learn the new task of generating leaf nodes from {\toc}?
\end{itemize}

\noindent Our key findings are as follows:
\begin{enumerate}
    \item {\sys} outperforms strong baselines such as BM25, DPR and fine-tuned DSI
    \item R@1 for Mistral is much lower than BM25 -- this suggests the LLM not only needs to learn the new task of picking the best section but also needs to ingest knowledge from the corpus
    \item DPR performs much better than BM25 but worse than DSI. This result is expected, as NV-Embed-v2 is not fine-tuned for {\ben}
    \item Tuning with {\toc} as input helps {\sys}: The only difference in the input between DSI and {\sys} is that {\sys} leverages {\toc}. {\sys} outperforms DSI by $7.4\%$, as it is able to access the entire structure of the book and better align it with the queries during fine-tuning. DSI on the other hand must learn the semantic alignment between queries and sections by looking at all the training data. We can thus conclude that \textbf{RQ1} and \textbf{RQ2} are answered in affirmative.
\end{enumerate}
%{\sys} achieves the best performance precisely because it jointly learns to generate from {\toc} and and ingesting corpus knowledge (Section ~\ref{subsec:train_task})
 % Table ~\ref{table:results_r1} demonstrates that Mistral performs poorly when used out-of-the-box to select the best section which can answer a user query achieving only a Recall@1 of 0.18. \todo[inline]{please check next two lines vineet}In contrast, BM25 performs significantly better than Mistral. 
 % This highlights the necessity of corpus knowledge ingestion. 
 
 % Table ~\ref{table:results_r1} indicates that the Mistral model, when used out-of-the-box to select the appropriate leaf node from the input {\toc}, does not yield optimal results, achieving an average Recall@1 of only 0.18. This suggests that section titles alone lack sufficient semantic information for the LLM to accurately match queries to the correct sections, highlighting the necessity of including both content and {\toc}. In contrast, BM25 performs significantly better than Mistral, as it builds the index using the text content from the book sections along with the training queries. This approach enriches the index with a diverse set of keywords derived from the original text and related queries.


% \fix{{\ben} is a challenging benchmark for embedding models, BM25 provides a competitive baseline better than DPR: Given that BM25 index is information rich (especially since it is indexed using training questions), this observation aligns with findings from the popular IR benchmark BEIR~\cite{thakur2021beir}, where BM25 consistently provided a competitive baseline, frequently outperforming early attempts with DPR trained on limited data ~\cite{thakur2021beir}. DPR finds it hard to align with the new task of title retrieval and diverse domain of \ben}


% {\sys} . It is worth highlighting that the same LLM, Mistral, was expertly tuned for both the DSI baseline and {\sys} using an identical set of queries, with the sole distinction being that DSI excluded the {\toc} from its input. This performance improvement of {\sys} over DSI clearly highlights the benefit of using {\toc}, as it allows {\sys} to access the entire structure of the book and better align it with the queries during fine-tuning. I This result answers \textbf{RQ1} affirmatively: using the {\toc} structure during LLM tuning enhances retriever performance. This also answers \textbf{RQ2} in positive indicating that an LLM can be taught the new task of generating the most relevant leaf node given a global corpus structure

% \textbf{{\sys} learns to generate from {\toc}}: The high Recall@1 answers \textbf{RQ2} in positive, indicating that with modern day LLMs it is indeed possible to reuse their inherent human instruction understanding capabilities to teach them a new task just by finetuning.
\eat{
\begin{table}[!htb]
\begin{tabular}{lllll}
\toprule
\textbf{Domain}    & \textbf{Total} & \textbf{(a) TOC!= DSI} & \textbf{(b) ToC-DSI} & \textbf{(c)DSI-ToC} \\
\midrule
Finance   & 16415 & 3883      & 1857    & 1053    \\
\midrule
Education & 15326 & 4387      & 2075    & 1014    \\
\midrule
Law       & 27472 & 9076      & 3764    & 2017    \\
\midrule
Medicine  & 17346 & 3355      & 1709    & 865     \\
\midrule
Nat\_sci  & 15719 & 4335      & 1860    & 1327    \\
\midrule
Soc\_sci  & 10971 & 2610      & 1209    & 487    \\
\bottomrule
\end{tabular}
\label{significance_testing}
\caption{Statistics on total number of instances per domain and number of instances where (a) DSI and ToC output differ (b) TOC is correct and DSI is incorrect (c) DSI is correct and ToC is incorrect  }
\end{table}
}
\textbf{Statistical significance testing} was done for the Recall@1 performance difference between DSI and {\sys}.  We follow the randomization test tailored to retrieval systems as described in ~\citet{sig_test}. The null hypothesis for the significance testing is that outputs from SystemA (DSI) and SystemB ({\sys}) may belong to the same underlying distribution and the difference of performance between SystemA (DSI) and SystemB ({\sys}) is because of the sampling variance. We test the null hypothesis with the traditional $p$ value of 0.05 at the domain level i.e. for each of the $6$ domains separately (by combining all the books from a single domain). The significance test results conclude that the difference in DSI and {\sys} is indeed statistically significant for all the $6$ domains where the null hypothesis is successfully invalidated. More specifically, through our sample runs for randomization test, we find that the probability of seeing this scale of Recall@1 performance difference as observed in Table~\ref{table:results_r1} when being randomly sampled from the same distribution falls below $p=0.05$. As with any statistical significance testing, the key takeaway from these experiments is a confirmation that the performance improvement seen in {\sys} is statistically significant and therefore, it is a property of the {\sys} design and not because of the data set size or data distribution in {\ben}.

\subsection{Ablations}
\label{subsec:ablation}
% We make two key conclusions by comparing {\sys} with DSI:
We conduct ablation studies to study the impact of using {\toc} as an input. Two key conclusions are as follows:

\noindent\textbf{{\toc} reduces hallucinations}: We define hallucination as the generation of a non-leaf node (invalid document identifier). Figure ~\ref{fig:ablation_a} illustrates the number of hallucinations per leaf node. We observed that the hallucination rate of {\sys} remains nearly constant at close to zero, regardless of the number of training examples. This indicates that {\sys} has effectively learned to generate outputs solely from the leaf nodes, making it inherently less prone to hallucination. In contrast, DSI, which does not use {\toc}, exhibits a higher hallucination rate than {\sys}, particularly for leaf nodes with fewer training examples.\\
\textbf{Remembering the structure is hard without {\toc}} Figure~\ref{fig:ablation_b} shows that for leaf nodes with less number of training examples, the Recall@1 difference between DSI and {\sys} is much higher. This supports our hypothesis, that without {\toc}, DSI needs to discover the book structure and also remember it within its parameters. DSI attempts to learn this via the training data on query and associated {\toc} leaf node pairs. Thus, leaf nodes which have low representations in training data are prone to be missed by DSI. In contrast, {\sys} does not need to ``remember'' the leaf nodes. As the number of training examples for a leaf node increase, Recall@1 gap between DSI and {\sys} decreases -- although {\sys} consistently achieves higher Recall@1 numbers.

\subsection{Error Analysis with Anecdotal Examples}
We observe that out-of-the-box Mistral's error rate is 86.20\%, and  26.81\% of its mistakes are as it predicts a non-leaf node, while 23.86\% are hallucinations. This suggests that Mistral out-of-box needs knowledge ingestion and task finetuning. For DSI, the error rate drops to 24.31\%, with 3.25\% of its predictions being non-leaf nodes (suggesting knowledge infusion helps). {\sys} has the lowest error rate of 18.67\%, with only 0.05\% of its predictions being non-leaf nodes.

\begin{figure}[h]
\caption{As number of training examples decrease, hallucination rate for DSI increases}
% \Description{DSI hallucinates as number of train examples decrease!}
\includegraphics[width=\linewidth]{images/hrate.png}
\label{fig:ablation_a}

\end{figure}

\begin{figure}[h]
% \caption{Recall@1 (averaged over a sliding window of 100) for all leaves across 18 books in our benchmark. Leaves are sorted by increasing number of training examples. {\sys} shows high accuracy even with a low number of training examples.}
\caption{Recall@1 for each leaf across {\ben}. {\sys} shows high accuracy even with a low number of training examples}
% \Description{{\sys} can learn from few training examples}
\includegraphics[width=\linewidth]{images/acc.png}
\label{fig:ablation_b}
\end{figure}


\eat{
\section{Significance tests}

\begin{table}[!htb]
\begin{tabular}{lllll}
\toprule
\textbf{Domain}    & \textbf{Total} & \textbf{(a) TOC!= DSI} & \textbf{(b) ToC-DSI} & \textbf{(c)DSI-ToC} \\
\midrule
Finance   & 16415 & 3883      & 1857    & 1053    \\
\midrule
Education & 15326 & 4387      & 2075    & 1014    \\
\midrule
Law       & 27472 & 9076      & 3764    & 2017    \\
\midrule
Medicine  & 17346 & 3355      & 1709    & 865     \\
\midrule
Nat\_sci  & 15719 & 4335      & 1860    & 1327    \\
\midrule
Soc\_sci  & 10971 & 2610      & 1209    & 487    \\
\bottomrule
\end{tabular}
\label{significance_testing}
\caption{Statistics on total number of instances per domain and number of instances where (a) DSI and ToC output differ (b) TOC is correct and DSI is incorrect (c) DSI is correct and ToC is incorrect  }
\end{table}
%\subsection{Anecdotal Examples}
%\label{sec:anecdotal_examples}
}
Table~\ref{table:pred_examples} lists a few qualitative examples. We discuss next the first example ``How do preschoolers react when caregivers and teachers belittle their autonomous actions'' from the book ``The Whole Child: Development in the Early Years'':

\textbf{Mistral} could correctly guess that the query can be answered by ``2.2.3 Initiative vs. Guilt (Preschool Years)'' using its world knowledge. It doesn't understand that it needs to pick only one option -- due to lack of the task training\\
\textbf{BM25} suffers from the known bias of keyword matching where and picks the section ``8.3.2. Moral Development'' purely because of highest lexical overlap with query tokens, which is the incorrect answer\\
\textbf{DPR} is used out of the box and does a word sensed matching where it predicts ``2.2.2 Autonomy vs. Shame/Doubt (Toddlerhood)'' being most relevant to the query tokens mentioning ``belittling of autonomous actions''. However, without the corpus specific training it falls short of identifying the correct section.\\
\textbf{DSI} is finetuned for the corpus but faulters to find the precise section. It errs just like DPR in this example but might have chosen a more sensible match ``2.2.2 Autonomy vs. Shame/Doubt (Preschool years)'' which takes into account the importance of keyword ``preschool'' for getting a better topic match. However, careful consideration shows this is a hallucinated header not present in the input {\toc}. This correlates with our ablation in figure~\ref{fig:ablation_a} on DSI hallucination.\\
\textbf{\sys}: Our proposed system {\sys} correctly identifies the gold leaf node as ``Initiative vs. Guilt (Preschool Years)''. 

The comparison between DPR, DSI and {\sys} is interesting as they all are built on top odf the same base LLM -- Mistral. The difference between DSI and {\sys} largely points towards the importance of having {\toc} as an input during training, which makes it easier for the LLM to align much better to the corpus. Other anecdotal examples in Table~\ref{table:pred_examples} also enumerate the same partial ordering seen among the baselines and provide more qualitative analysis to support the claim that the design of {\sys} helps it do a much better retrieval.

% We show some anecdotal examples from {\ben} in Table~\ref{table:pred_examples} which further demonstrates the above analysis with real data instances. For example1, we can see how Mistral errs in producing a leaf node showing the lack of task training to generate from {\toc} leaf nodes. Example 2 and example 3 show two different cases where DSI gets it incorrect for different reasons. 

% For example 2, the gold leaf node had only $4$ training examples compared to an average of $18.64$ training examples across all leaf nodes from the same book. As we have already seen in Figure~\ref{fig:ablation_b}, DSI finds it hard to remember the gold leaf node \textit{19.3 Oxidation of Aldehydes and Ketones} possibly due to lack of enough training examples and goes on to produce a hallucinated leaf node.

% For example 3 too,  DSI errs because lack of enough training examples for the gold leaf node \textit{9.5.1 What are non-democracies?}, which is $2$ as compared to an average of $11.79$ training examples for the book. However, DSI instead of hallucinating, predicts a different leaf node \textit{6.2.1 The Experimental Method } which is seemingly relevant and has $21$ training examples for DSI to remember it. 

% For all 3 examples {\sys} produces the right leaf node by learning to utilize the input {\toc} effectively.






% \begin{table}[ht]
% \caption{Predictions for the book \href{https://open.umn.edu/opentextbooks/textbooks/221}{Teaching in a Digital Age}}
% \label{table:pred_examples}
% \begin{center}
% % \addtolength{\tabcolsep}{-0.4em}
% \begin{tabular}{lrl}
% \hline
% \multirow{4}{*}{Ex1}&\textbf{Query}:& How do cultural factors influence discussions in online education? \\
% & \textbf{Gold}:&4.4 Online collaborative learning \\
% & \textbf{Mistral}:&9.10 Deciding (Chapter 9) \\
% & \textbf{DSI}:&12.10 Step eight: communicate, communicate, communicate \\
% & \textbf{{\sys}}:&4.4 Online collaborative learning \\
% \hline
% \multirow{5}{*}{Ex2}&\textbf{Query}:& What measures must be taken to guarantee a gradual advancement \\
% & & in a student's ability over the course of a curriculum?\\
% & \textbf{Gold}:&6.5 Developing skills\\
% & \textbf{Mistral}:&12.6 Step six: set appropriate learning goals.\\
% & \textbf{DSI}:&8.5 Computing\\
% & \textbf{{\sys}}:&6.5 Developing skills\\
% \hline
% \multirow{5}{*}{Ex3}&\textbf{Query}:&What is the suggested time allocation for students studying \\
% & & a course or program?\\
% & \textbf{Gold}:&12.9 Step seven: design course structure and learning activities\\
% & \textbf{Mistral}:&12.6 Step six: set appropriate learning goals\\
% & \textbf{DSI}:&9.3 Ease of Use	\\
% & \textbf{{\sys}}:&12.9 Step seven: design course structure and learning activities\\
% \end{tabular}
% \end{center}
% \end{table}

% We look at some predictions for the book ``Teaching in a Digital Age`` in Table ~\ref{table:pred_examples}. Table of contents for this book is highlighted in Figure~\ref{fig:edu_book3_toc}. In Ex1, Mistral generates incorrect leaf as output and it indicates that it does not understand that it only needs to generate the leaf node, even though it is clearly instructed to do so in the prompt (Section ~\ref{subsec:toc_prompt}). DSI system makes an incorrect prediction as it is not able to leverage the fact that a query about online education could probably be answered by Chapter 4 which is titled ``Methods for teaching with an online focus``. {\sys} makes the correct prediction! In Ex2, DSI outputs a completely unrelated section, as it is unable to leverage the rich information present in Table of Contents -- that ``Chapter 6'' talks about learning and thus ``6.5'' is a better fit for the user query. Ex3 highlights another issue where Mistral model predicts incorrect (hallucinates) leaf node and DSI model predicts a completely unrelated leaf node.

